\section{Reactive: WebSocket \& SSE}
\label{sec:reactive}

Marking an action \texttt{method: reactive} tells Emi the endpoint is a
WebSocket, full-duplex channel. The JS compiler generates a
\texttt{WebSocketX} class --- the base every reactive action's generated
code sits on top of --- which extends the standard \texttt{WebSocket}
with the exception that \texttt{message} and \texttt{send} are overridden
to accept generic types, and the constructor takes a factory for message
creation when working with JSON messages. Other generators (React, for
instance) wrap the instance rather than replacing it, layering a
\texttt{useWebSocketX} hook on top.

\begin{lstlisting}
actions:
  - name: userStream
    url: ws://localhost:8081
    method: reactive
    in:
      fields:
        - {name: min,   type: int}
        - {name: max,   type: int}
        - {name: count, type: int}
    out:
      fields:
        - {name: number, type: int}
\end{lstlisting}

produces a complete typed request/response pair plus an action wrapper:

\begin{lstlisting}
export class UserStreamAction {
  static URL = "ws://localhost:8081";
  static Method = "REACTIVE";
  static Create = (overrideUrl?, qs?, options) => {
    const Cls = options?.SocketClass ??
      WebSocketX<UserStreamActionReq, UserStreamActionRes>;
    return new Cls(url, undefined, {
      MessageFactoryClass: UserStreamActionRes,
    });
  };
}
\end{lstlisting}

The generated request (\texttt{UserStreamActionReq}) and response
(\texttt{UserStreamActionRes}) classes are ordinary Emi DTO classes ---
same private fields, typed getters/setters, and auto-init guarantees as
any other generated class (\S\ref{sec:js-compiler}); the only thing
\texttt{reactive} changes is the transport they travel over. The
generated JavaScript works in both Node.js and browser environments,
though the client side is the primary focus.

\begin{pullquote}
SSE gets the same treatment as WebSocket in spirit: a plain \texttt{get}
action whose response streams chunked messages compiles to typed
messages on the client, and \texttt{--tags react} adds a matching
\texttt{useSse} hook right alongside \texttt{useWebSocketX}.
\end{pullquote}
